\chapter{Commuting-Projector Models and the Toric Code}
\label{chap:commuting-projector-toric-code}

The previous part studied interacting one-dimensional quantum chains whose exact
solvability comes from factorized scattering and the Yang--Baxter equation.  In
this part we move to two spatial dimensions.  The first solvable structure is
very different from Bethe-ansatz integrability: the Hamiltonian is a sum of
local terms that commute with one another.  Such models are exactly solvable
because all local conserved projectors can be diagonalized simultaneously.

The toric code is the canonical example.  It is simultaneously
\begin{enumerate}[label=(\roman*)]
    \item a stabilizer Hamiltonian;
    \item a commuting-projector model;
    \item a lattice \(\ZZ_2\) gauge theory at zero gauge coupling;
    \item the simplest fixed-point model of intrinsic topological order.
\end{enumerate}
It is therefore an ideal first model for two-dimensional quantum exact
solvability.  Its solution does not require Fourier transforms, Bogoliubov
rotations, or Bethe rapidities.  Instead, the spectrum follows from a finite set
of mutually commuting \(\ZZ_2\)-valued operators, while its physical content is
encoded in nonlocal string operators.

\section{Commuting-projector solvability}
\label{sec:commuting-projector-principle}

Let \(\calH\) be a many-body Hilbert space and suppose that a Hamiltonian can be
written as
\begin{equation}
\label{eq:commuting-projector-general-form}
    H=\sum_{\alpha} E_{\alpha} P_{\alpha},
    \qquad
    P_{\alpha}^2=P_{\alpha},
    \qquad
    [P_{\alpha},P_{\beta}]=0 .
\end{equation}
Then all \(P_{\alpha}\) can be diagonalized simultaneously.  The many-body
spectrum is reduced to the problem of finding which binary eigenvalue patterns
\(p_{\alpha}\in\{0,1\}\) are compatible with the algebraic constraints among the
projectors.

Equivalently, if each local term is written in terms of an involution
\(Q_{\alpha}^2=\id\),
\begin{equation}
\label{eq:involution-projector-relation}
    P_{\alpha}=\frac{1-Q_{\alpha}}{2},
    \qquad
    Q_{\alpha}=\id-2P_{\alpha},
\end{equation}
then exact solvability is the statement that the commuting stabilizers
\(Q_{\alpha}\) admit simultaneous eigenvalues \(q_{\alpha}=\pm1\).

\begin{definition}[Commuting-projector Hamiltonian]
A local Hamiltonian is called a commuting-projector Hamiltonian if it can be
written as a sum of local projectors satisfying \cref{eq:commuting-projector-general-form}.
If the ground state minimizes every local projector separately, the Hamiltonian
is called frustration-free.
\end{definition}

\begin{remark}[Exact solvability versus integrability]
A commuting-projector model has an extensive set of strictly local conserved
operators.  This is stronger than ordinary quantum integrability in one
dimension, where conserved charges are usually quasilocal or generated by a
transfer matrix.  However, the strength of the solvability is also more
specialized: a generic perturbation that destroys commutativity usually cannot be
solved by the same method.  The toric code should therefore be viewed as an
exactly solvable fixed point, not as a representative of a continuously soluble
family in the Bethe-ansatz sense.
\end{remark}

The simplest nontrivial commuting-projector models are stabilizer Hamiltonians.
There one has \(N\) qubits and an Abelian subgroup \(\mathcal{S}\) of the
Pauli group, generated by mutually commuting Pauli strings \(S_a\).  The common
\(+1\) eigenspace
\begin{equation}
\label{eq:stabilizer-code-space}
    \calH_{\rm code}
    =
    \{\ket{\psi}\in(\CC^2)^{\otimes N}:S_a\ket{\psi}=\ket{\psi}
    \text{ for all }a\}
\end{equation}
is the ground space of a positive-semidefinite Hamiltonian
\begin{equation}
\label{eq:stabilizer-positive-hamiltonian}
    H_{\rm stab}=\sum_a \frac{1-S_a}{2} .
\end{equation}
If \(r\) of the stabilizer generators are independent, then
\begin{equation}
\label{eq:stabilizer-dimension-counting}
    \dim\calH_{\rm code}=2^{N-r} .
\end{equation}
This counting formula will give the toric-code ground-state degeneracy on a
closed surface.

\section{Lattice, Hilbert space, and Hamiltonian}
\label{sec:toric-code-hamiltonian}

Consider an oriented square lattice embedded on a closed orientable surface
\(\Sigma_g\) of genus \(g\).  The set of vertices, edges, and plaquettes is
denoted by
\begin{equation}
\label{eq:cellulation-sets}
    V,
    \qquad
    E,
    \qquad
    P,
\end{equation}
with cardinalities \(N_v=\abs{V}\), \(N_e=\abs{E}\), and \(N_p=\abs{P}\).
A spin-\(1/2\) degree of freedom lives on every edge \(\ell\in E\).  Thus
\begin{equation}
\label{eq:toric-code-hilbert-space}
    \calH_{\rm TC}
    =\bigotimes_{\ell\in E}\CC^2_{\ell},
    \qquad
    \dim\calH_{\rm TC}=2^{N_e} .
\end{equation}
We write Pauli operators on edge \(\ell\) as \(\sigma^x_{\ell}\),
\(\sigma^y_{\ell}\), and \(\sigma^z_{\ell}\).

For each vertex \(s\in V\), define the star operator
\begin{equation}
\label{eq:toric-code-star-operator}
    A_s=
    \prod_{\ell\ni s}\sigma^x_{\ell},
\end{equation}
where the product is over the edges incident on \(s\).  For each plaquette
\(p\in P\), define the plaquette operator
\begin{equation}
\label{eq:toric-code-plaquette-operator}
    B_p=
    \prod_{\ell\in\partial p}\sigma^z_{\ell},
\end{equation}
where the product is over the boundary edges of \(p\).  On a square lattice each
star and plaquette contains four edges, but the algebraic construction only
requires that the lattice be a cellulation of a closed oriented surface.

The toric-code Hamiltonian is
\begin{equation}
\label{eq:toric-code-hamiltonian}
    H_{\rm TC}
    =
    -J_e\sum_{s\in V}A_s
    -J_m\sum_{p\in P}B_p,
    \qquad
    J_e,J_m>0 .
\end{equation}
Equivalently, up to an additive constant,
\begin{equation}
\label{eq:toric-code-projector-form}
    H_{\rm TC}
    =
    2J_e\sum_{s\in V}\frac{1-A_s}{2}
    +2J_m\sum_{p\in P}\frac{1-B_p}{2}
    -(J_eN_v+J_mN_p).
\end{equation}
The terms \((1-A_s)/2\) and \((1-B_p)/2\) are projectors onto violated star and
plaquette constraints.

\begin{remark}[Naming convention]
The notation \(J_e\) and \(J_m\) anticipates the excitation content.  A violated
star, \(A_s=-1\), is called an electric charge or \(e\) particle.  A violated
plaquette, \(B_p=-1\), is called a magnetic flux or \(m\) particle.  The model
is self-dual under exchange of direct and dual lattices together with
\(A_s\leftrightarrow B_p\), \(J_e\leftrightarrow J_m\), and
\(\sigma^x\leftrightarrow\sigma^z\).
\end{remark}

\section{Commutativity and simultaneous diagonalization}
\label{sec:toric-code-commutativity}

The toric code is exactly solvable because every star and every plaquette term
commutes.

\begin{proposition}[Toric-code stabilizer algebra]
The operators \(A_s\) and \(B_p\) satisfy
\begin{equation}
\label{eq:toric-code-stabilizer-algebra}
    A_s^2=B_p^2=\id,
    \qquad
    [A_s,A_{s'}]=0,
    \qquad
    [B_p,B_{p'}]=0,
    \qquad
    [A_s,B_p]=0 .
\end{equation}
\end{proposition}

\begin{proof}
The first identity follows from \((\sigma^x)^2=(\sigma^z)^2=\id\).  Two star
operators contain only \(\sigma^x\) operators, hence commute.  Two plaquette
operators contain only \(\sigma^z\) operators, hence commute.

It remains to check \([A_s,B_p]=0\).  The operators \(A_s\) and \(B_p\) can only
fail to commute on edges shared by the star \(s\) and the plaquette \(p\).  On a
shared edge,
\begin{equation}
\label{eq:pauli-anticommute-local}
    \sigma^x_{\ell}\sigma^z_{\ell}
    =
    -\sigma^z_{\ell}\sigma^x_{\ell} .
\end{equation}
On a square lattice embedded without boundary, a star and a plaquette share
either zero or two edges.  Therefore moving \(A_s\) through \(B_p\) produces the
factor \((-1)^n\), where \(n\in\{0,2\}\).  Since \((-1)^n=+1\), the two
operators commute.
\end{proof}

Because of \cref{eq:toric-code-stabilizer-algebra}, the Hamiltonian can be
diagonalized by simultaneous eigenstates
\begin{equation}
\label{eq:toric-code-simultaneous-eigenvalues}
    A_s\ket{\psi}=a_s\ket{\psi},
    \qquad
    B_p\ket{\psi}=b_p\ket{\psi},
    \qquad
    a_s,b_p\in\{+1,-1\} .
\end{equation}
The energy of such a state is
\begin{equation}
\label{eq:toric-code-energy-eigenvalue}
    E[\{a_s\},\{b_p\}]
    =
    -J_e\sum_{s\in V}a_s
    -J_m\sum_{p\in P}b_p .
\end{equation}
Thus the full spectrum is determined once we know the allowed stabilizer
eigenvalue patterns and their degeneracies.

On a closed surface there are two global relations:
\begin{equation}
\label{eq:toric-code-global-relations}
    \prod_{s\in V} A_s=\id,
    \qquad
    \prod_{p\in P} B_p=\id .
\end{equation}
The first relation holds because every edge touches exactly two vertices and
therefore appears twice in \(\prod_s A_s\).  The second holds because every edge
belongs to exactly two plaquette boundaries and therefore appears twice in
\(\prod_p B_p\).

It follows that the number of star violations and the number of plaquette
violations must both be even:
\begin{equation}
\label{eq:toric-code-even-parity-constraints}
    \prod_{s\in V}a_s=1,
    \qquad
    \prod_{p\in P}b_p=1 .
\end{equation}
In particular, on a closed surface a single isolated \(e\) or \(m\) excitation
cannot be created by a finite local operator.  Excitations are created in pairs,
or else one must use open boundaries or punctures.

\section{Ground states as flat \texorpdfstring{\(\ZZ_2\)}{Z2} gauge fields}
\label{sec:toric-code-ground-states}

The ground-state conditions are
\begin{equation}
\label{eq:toric-code-ground-state-conditions}
    A_s\ket{\Omega}=\ket{\Omega},
    \qquad
    B_p\ket{\Omega}=\ket{\Omega}
\end{equation}
for all vertices \(s\) and plaquettes \(p\).  These equations have a useful gauge
theory interpretation.  The \(\sigma^z\)-basis assigns a \(\ZZ_2\) variable to
each edge.  The plaquette constraint \(B_p=+1\) says that the \(\ZZ_2\) flux
through every plaquette is trivial.  The star operator \(A_s\) flips all edge
variables adjacent to a vertex; it acts as a local \(\ZZ_2\) gauge
transformation.  Thus the toric-code ground states are gauge-invariant
superpositions of flat \(\ZZ_2\) connections.

A concrete construction is obtained as follows.  Let \(
\ket{0_z}\) be the
state with \(\sigma^z_{\ell}=+1\) for all edges.  It already satisfies
\(B_p\ket{0_z}=\ket{0_z}\) for every plaquette.  The projector onto the
\(+1\)-eigenspace of all star operators is
\begin{equation}
\label{eq:toric-code-star-projector}
    \Pi_A
    =
    \prod_{s\in V}\frac{1+A_s}{2} .
\end{equation}
Therefore
\begin{equation}
\label{eq:toric-code-plane-ground-state}
    \ket{\Omega_0}
    \propto
    \Pi_A\ket{0_z}
\end{equation}
is a ground state in the topologically trivial sector.  In the \(\sigma^z\)
basis, this state is an equal-amplitude superposition of closed-loop
configurations.  The word ``closed'' means that every vertex is touched by an
even number of flipped edges; this is precisely the no-charge condition
\(A_s=+1\).

On a torus there are additional sectors distinguished by noncontractible
\(\ZZ_2\) fluxes.  Local star operations cannot change these global fluxes, so
one obtains four independent ground states rather than one.  We derive this
counting in \cref{sec:toric-code-topological-degeneracy}.

\section{String operators and anyon excitations}
\label{sec:toric-code-anyons}

The elementary excitations are defects of the stabilizer constraints.

\begin{definition}[Electric and magnetic excitations]
A vertex \(s\) with \(A_s=-1\) is an electric charge, denoted \(e\).  A plaquette
\(p\) with \(B_p=-1\) is a magnetic flux, denoted \(m\).  The composite excitation
\(\epsilon=e\times m\) is obtained by placing an \(e\) and an \(m\) at the same
position in the continuum topological description.
\end{definition}

The energy cost of one violated star is \(2J_e\), because its contribution to
\cref{eq:toric-code-hamiltonian} changes from \(-J_e\) to \(+J_e\).  Similarly,
one violated plaquette costs \(2J_m\).  On a closed surface the minimal locally
created excitations are pairs, with energies
\begin{equation}
\label{eq:toric-code-pair-energy-costs}
    \Delta E_{e\text{-pair}}=4J_e,
    \qquad
    \Delta E_{m\text{-pair}}=4J_m .
\end{equation}

\subsection{Electric strings}

Let \(\gamma\) be a path on the direct lattice, i.e. a sequence of edges.  Define the electric string operator
\begin{equation}
\label{eq:toric-code-electric-string}
    W_e(\gamma)
    =
    \prod_{\ell\in\gamma}\sigma^z_{\ell} .
\end{equation}
Since \(W_e(\gamma)\) is made of \(\sigma^z\) operators, it commutes with all
plaquette operators \(B_p\).  It anticommutes with a star operator \(A_s\) if and
only if \(s\) is an endpoint of \(\gamma\).  Therefore an open electric string
creates or moves a pair of \(e\) excitations at its endpoints:
\begin{equation}
\label{eq:e-string-endpoint-action}
    A_s W_e(\gamma)
    =
    \begin{cases}
        -W_e(\gamma)A_s, & s\in\partial\gamma,\\
        \phantom{-}W_e(\gamma)A_s, & s\notin\partial\gamma .
    \end{cases}
\end{equation}
For a closed path, \(\partial\gamma=\varnothing\), the string operator commutes
with the Hamiltonian.

\subsection{Magnetic strings}

Let \(\gamma^\ast\) be a path on the dual lattice.  It crosses a set of direct
lattice edges.  Define
\begin{equation}
\label{eq:toric-code-magnetic-string}
    W_m(\gamma^\ast)
    =
    \prod_{\ell\perp\gamma^\ast}\sigma^x_{\ell} .
\end{equation}
This operator commutes with all star operators \(A_s\), but it anticommutes with
the two plaquette operators at the endpoints of \(\gamma^\ast\).  Thus an open
magnetic string creates or moves a pair of \(m\) fluxes.

\subsection{Fusion rules}

Because \(e\) and \(m\) excitations are \(\ZZ_2\)-valued defects, bringing two
identical excitations together annihilates them.  The fusion rules are
\begin{equation}
\label{eq:toric-code-fusion-rules}
    e\times e=1,
    \qquad
    m\times m=1,
    \qquad
    e\times m=\epsilon,
    \qquad
    \epsilon\times\epsilon=1 .
\end{equation}
The set of anyon types is therefore
\begin{equation}
\label{eq:toric-code-anyon-types}
    \{1,e,m,\epsilon\}\cong \ZZ_2\times\ZZ_2
\end{equation}
under fusion.

\section{Mutual statistics}
\label{sec:toric-code-mutual-statistics}

The nontrivial topological content of the toric code is not the energy of the
excitations, which is completely local, but their braiding statistics.  Let
\(\gamma\) be a direct-lattice path and \(\gamma^\ast\) a dual-lattice path.  The
corresponding string operators obey
\begin{equation}
\label{eq:toric-code-string-crossing-algebra}
    W_e(\gamma)W_m(\gamma^\ast)
    =
    (-1)^{I(\gamma,\gamma^\ast)}
    W_m(\gamma^\ast)W_e(\gamma),
\end{equation}
where \(I(\gamma,\gamma^\ast)\) is the number of intersections modulo \(2\).
The reason is local: at each crossing the two strings act on the same edge by
\(\sigma^z\) and \(\sigma^x\), which anticommute.

If an \(e\) particle is adiabatically moved once around an \(m\) particle, the
worldline process is represented algebraically by a closed electric string that
intersects a magnetic string once.  Hence the wave function acquires the phase
\begin{equation}
\label{eq:toric-code-mutual-minus-one}
    \theta_{em}=-1 .
\end{equation}
Thus \(e\) and \(m\) are mutual semions.

The individual self-statistics are bosonic,
\begin{equation}
\label{eq:toric-code-self-statistics}
    \theta_e=+1,
    \qquad
    \theta_m=+1,
\end{equation}
because exchanging two identical \(e\)'s or two identical \(m\)'s can be
represented without producing a nontrivial string crossing.  Their composite
\(\epsilon=e\times m\) has fermionic self-statistics,
\begin{equation}
\label{eq:toric-code-epsilon-fermion}
    \theta_{\epsilon}=-1 .
\end{equation}
This is one of the simplest examples where a purely spin Hamiltonian supports an
emergent fermionic quasiparticle.

\begin{remark}[Why this is topological]
The phase in \cref{eq:toric-code-string-crossing-algebra} depends only on the
intersection number modulo \(2\), not on the metric length or detailed shape of
the paths.  This is the operational meaning of topological statistics in the
toric-code fixed-point model.
\end{remark}

\section{Topological degeneracy on closed surfaces}
\label{sec:toric-code-topological-degeneracy}

We now count the ground-state degeneracy on a closed orientable surface
\(\Sigma_g\).  The torus corresponds to \(g=1\).

The toric code has \(N_e\) physical qubits.  The stabilizer group is generated by
\(N_v\) star operators and \(N_p\) plaquette operators, but the two relations in
\cref{eq:toric-code-global-relations} remove two independent generators.  Hence
the number of independent stabilizers is
\begin{equation}
\label{eq:toric-code-independent-stabilizers}
    r=N_v+N_p-2 .
\end{equation}
By the stabilizer counting formula \cref{eq:stabilizer-dimension-counting},
\begin{equation}
\label{eq:toric-code-gsd-counting-first}
    \dim\calH_{\rm GS}
    =
    2^{N_e-r}
    =
    2^{N_e-N_v-N_p+2} .
\end{equation}
Using the Euler characteristic
\begin{equation}
\label{eq:euler-characteristic-genus}
    N_v-N_e+N_p=2-2g,
\end{equation}
we obtain
\begin{equation}
\label{eq:toric-code-gsd-genus-g}
    \boxed{
    \dim\calH_{\rm GS}(\Sigma_g)=2^{2g} .
    }
\end{equation}
On a torus, \(g=1\), this gives
\begin{equation}
\label{eq:toric-code-gsd-torus}
    \dim\calH_{\rm GS}(T^2)=4 .
\end{equation}

\begin{theorem}[Toric-code topological degeneracy]
On a closed orientable surface of genus \(g\), the toric-code Hamiltonian has
\(2^{2g}\) exactly degenerate ground states.  In particular, on the torus the
ground-state degeneracy is four.
\end{theorem}

\begin{remark}[Why the degeneracy is not symmetry breaking]
The degeneracy in \cref{eq:toric-code-gsd-genus-g} is not caused by a local order
parameter.  It is tied to noncontractible cycles of the spatial manifold.  On a
sphere, \(g=0\), the ground state is unique.  On a torus, the four ground states
cannot be distinguished by any strictly local operator in the thermodynamic
limit; they are distinguished by nonlocal Wilson loop operators.
\end{remark}

\section{Logical Wilson loops on the torus}
\label{sec:toric-code-logical-loops}

Let \(\alpha\) and \(\beta\) be the two noncontractible cycles of a torus.  Choose
direct-lattice cycles \(\gamma_{\alpha}\), \(\gamma_{\beta}\) and dual-lattice
cycles \(\gamma^\ast_{\alpha}\), \(\gamma^\ast_{\beta}\).  The corresponding
closed string operators are
\begin{equation}
\label{eq:toric-code-logical-loops}
    \overline{Z}_{\alpha}=W_e(\gamma_{\alpha}),
    \qquad
    \overline{Z}_{\beta}=W_e(\gamma_{\beta}),
    \qquad
    \overline{X}_{\alpha}=W_m(\gamma^\ast_{\alpha}),
    \qquad
    \overline{X}_{\beta}=W_m(\gamma^\ast_{\beta}) .
\end{equation}
All four operators commute with \(H_{\rm TC}\), because their strings are closed.
However, they do not all commute with one another.  Their algebra is determined
by intersection numbers:
\begin{equation}
\label{eq:toric-code-logical-loop-algebra}
    \overline{Z}_{\alpha}\overline{X}_{\beta}
    =
    -\overline{X}_{\beta}\overline{Z}_{\alpha},
    \qquad
    \overline{Z}_{\beta}\overline{X}_{\alpha}
    =
    -\overline{X}_{\alpha}\overline{Z}_{\beta},
\end{equation}
while pairs with zero intersection commute.  Thus the ground-state space on the
torus carries two encoded qubits.  One may choose \(\overline{Z}_{\alpha}\) and
\(\overline{Z}_{\beta}\) as a commuting set of labels, while the corresponding
\(\overline{X}\)-loops act as logical bit flips.

This is the stabilizer-code interpretation of the toric code.  Local errors
create anyon pairs, while noncontractible error strings implement logical
operators.  The code distance is set by the shortest noncontractible cycle of
the lattice.

\section{Local indistinguishability and robustness}
\label{sec:toric-code-local-indistinguishability}

Let \(\ket{\Omega_a}\) and \(\ket{\Omega_b}\) be two different ground states on a
large torus, distinguished by Wilson-loop eigenvalues.  A local operator
\(O_{\rm loc}\) supported on a contractible region cannot measure a
noncontractible flux.  Within the exactly solvable fixed-point model, this gives
the selection rule
\begin{equation}
\label{eq:toric-code-local-indistinguishability}
    \bra{\Omega_a}O_{\rm loc}\ket{\Omega_b}
    =
    C(O_{\rm loc})\delta_{ab} .
\end{equation}
This holds for operators whose support is small compared with the system size, up to
finite-size corrections when perturbations are added.

Physically, this is why the toric-code degeneracy is stable against weak local
perturbations in the thermodynamic limit.  A perturbation can split the ground
states only through virtual processes that create an anyon pair, move one anyon
around a noncontractible cycle, and annihilate the pair.  The amplitude of such a
process is exponentially small in the linear system size in the gapped phase.

\begin{principle}[Exact fixed point of topological order]
The toric code is not merely a solvable Hamiltonian.  It is a fixed-point
representative of a stable gapped phase.  Its universal data are the anyon types
\(\{1,e,m,\epsilon\}\), their fusion rules, their mutual statistics, and the
genus-dependent ground-state degeneracy \(2^{2g}\).
\end{principle}

\section{Summary and relation to later chapters}
\label{sec:toric-code-summary}

The exact solution of the toric code is based on four elementary facts:
\begin{enumerate}[label=(\roman*)]
    \item the Hamiltonian is a sum of commuting stabilizers;
    \item star and plaquette violations define \(e\) and \(m\) excitations;
    \item open strings create anyon pairs, while closed noncontractible strings
    act within the ground-state space;
    \item the ground-state degeneracy on \(\Sigma_g\) is \(2^{2g}\).
\end{enumerate}
This gives the simplest example of intrinsic topological order in these notes.
The following chapters generalize the same structural ideas in two directions.
Quantum double models replace the group \(\ZZ_2\) by a finite group \(G\), leading
in general to non-Abelian anyons.  The Kitaev honeycomb model keeps spin-\(1/2\)
degrees of freedom in two dimensions but solves them by mapping to Majorana
fermions moving in a static \(\ZZ_2\) gauge field.  Thus the toric code is the
commuting-projector endpoint of a broader family of exactly solvable
spin-liquid models.
